\section{经典统计力学}
\subsection{从拉格朗日力学到相空间}
在\href{https://github.com/Astolfo-Official/Baby-Math-on-Quantum-Chemistry/blob/main/%E3%80%8A%E9%87%8F%E5%AD%90%E5%8C%96%E5%AD%A6%E4%B8%AD%E7%9A%84%E6%95%B0%E5%AD%A6%E3%80%8B2.0.pdf}{我搬运的一本速通教材的第四章}中，我们简单过了一遍经典力学。感兴趣的可以先阅读一下前置章节，这里我们会跳过许多细节直接叙述经典力学的两大形式后进入经典相空间理论的介绍。经典相空间理论也是经典力学的一种等价表述方式，但是在处理大量自由度问题也即经典统计力学中，相空间理论有其他三种形式不可比拟的优势。简单来说，在经典统计力学中，我们最后要处理的对象并不是某一条具体轨道，而是\textbf{大量可能微观状态在相空间中的分布}。这对于牛顿力学，拉格朗日力学和哈密顿力学来说都是一个全新的挑战。

在拉格朗日力学中，设体系有$3N$个广义坐标$q_i$，其运动由拉格朗日量$L(\bm{q},\dot{\bm{q}},t)=T(\bm{q},\dot{\bm{q}},t)-V(\bm{q},\dot{\bm{q}},t)$决定，其中$T$是动能，$V$是势能。对应的运动学方程即为 Euler-Lagrange 方程：
\[\dv{}{t}\left(\pdv{L}{\dot{q}_i}\right)-\pdv{L}{q_i}=0,\quad (i=1,\ldots,3N)\]
容易验证上述方程等价于牛顿第二定律。进一步地，我们还定义与广义坐标$q_i$对应的广义动量：
\[p_i=\pdv{L}{\dot{q}_i}\]
这个关系可以反解出$\dot{q}_i=\dot{q}_i(\bm{q},\bm{p},t)$，则我们可以对拉格朗日量做 Legendre 变换 [章节\ref{sec:legendre}]，进而定义哈密顿量$H(\bm{q},\bm{p},t)=\bm{p}\dot{\bm{q}}-L(\bm{q},\dot{\bm{q}},t)$。如果势能$V$与体系动量（或速度）无关，哈密顿量就是我们熟悉的总能量。把$H$看成$q_i,p_i$的函数后，对其微分并逐项比较可得哈密顿正则方程：
\[\dot{q}_i=\pdv{H}{p_i},\quad \dot{p}_i=-\pdv{H}{q_i}\]

在哈密顿形式下体系在某一时刻$t$的微观状态由全部坐标和动量唯一确定$\xi=(q_1,\ldots,q_f,p_1,\ldots,p_f)$，所有这样的点构成一个$6N$维空间，称为\textbf{相空间}；相空间中的一个点称为相点，它完整代表一个微观状态。体系随时间演化时，相点便在相空间中描出一条轨线，称为\textbf{相轨线}。如果\textcolor{blue}{\textbf{哈密顿量不显示含时间}}，则能量守恒，不同轨线被限制在各自能量面上，相空间中同一点出发的轨迹全部等同，任意交叉的轨迹都代表同一个微观状态。

至此我们把研究对象从具体的轨道转向了相空间中的分布。这时就可以引入统计的概念：对同一个宏观系统，在给定时刻所有可能微观状态组成一个系综；若记相空间概率密度为$\rho(\bm{q},\bm{p},t)$，则满足归一化：
\[\int\rho(\bm{q},\bm{p},t)\dd\bm{\Gamma}=\frac{1}{h^{3N}}\int\rho(\bm{q},\bm{p},t)\dd^{3N}\bm{q}\dd^{3N}\bm{p}=1\]
这里$h$是相空间微分$\dd{\bm{\Gamma}}$的体积，一般取为 Planck 常数。于是任意可观测量$F(\bm{p},\bm{q},t)$的系综平均都可写为：
\[\langle F \rangle=\int F(\bm{p},\bm{q},t)\rho(\bm{q},\bm{p},t)\dd{\bm{\Gamma}}\]
考虑对$F$求对时间$t$的全导数并带入哈密顿正则方程：
\[\dv{F}{t}=\pdv{F}{t}+\sum_i\left(\pdv{F}{q_i}\dot{q}_i+\pdv{F}{p_i}\dot{p}_i\right)=\pdv{F}{t}+\sum_i\left(\pdv{F}{q_i}\pdv{H}{p_i}-\pdv{F}{p_i}\pdv{H}{q_i}\right)=\pdv{F}{t}+\{F,H\}\]
其中 Poisson 括号定义为：
\[\{F,G\}=\sum_i\left(\pdv{F}{q_i}\pdv{G}{p_i}-\pdv{F}{p_i}\pdv{G}{q_i}\right)\]
另一方面，概率守恒要求概率密度$\rho(\bm{q},\bm{p},t)$满足连续性方程：
\[\begin{aligned}
\pdv{\rho}{t}+\nabla\cdot(\rho \dot{\xi})&=\pdv{\rho}{t}+\nabla\rho\cdot\dot{\xi}+\rho(\nabla\cdot\dot{\xi})=\pdv{\rho}{t}+\nabla\rho\cdot\dot{\xi}+\rho\sum_{i=1}^{f}\left(\pdv{\dot{q}_i}{q_i}+\pdv{\dot{p}_i}{p_i}\right)\\
&=\pdv{\rho}{t}+\nabla\rho\cdot\dot{\xi}+\rho\sum_{i=1}^{f}\left(\pdv{}{q_i}\pdv{H}{p_i}-\pdv{}{p_i}\pdv{H}{q_i}\right)=\pdv{\rho}{t}+\nabla\rho\cdot\dot{\xi}=0 \quad \Leftrightarrow \quad \dv{\rho}{t}=0
\end{aligned}\]
这就是 Liouville 方程。它告诉我们：\textcolor{blue}{\textbf{在精确的哈密顿演化下，相空间中的概率流像不可压缩流体一样运动}}，局部概率密度沿着相轨线保持不变。若哈密顿量不显含时间，则能量守恒，相点始终停留在某个能量面，记为$\Gamma_E:\ H(\bm{q},\bm{p})=E$，在 Liouville 方程的约束下，轨迹的运动可以与流类比，我们考虑两种流：
\begin{itemize}
\item 1. 非遍历性流（Non-ergodic）：轨线在能量面上只占据一个子区域，如周期性运动；
\item 2. 遍历性流（Ergodic）：轨线在能量面上占据整个区域。
\end{itemize}
如果体系的流是遍历的，那么一条足够长的轨线在长时间平均意义下就能代表整个能量面的平均，即：
\[\overline{f}=\lim_{T\to\infty}\frac{1}{T}\int_0^T f(\bm{q},\bm{p})\dd t=\frac{1}{\Sigma_E}\int f(\bm{q},\bm{p})\delta(H(\bm{q},\bm{p})-E)\dd{\bm{\Gamma}}=\langle f \rangle_E\]
此时可定义能量面截面（微正则配分函数）与微正则分布：
\[\Sigma_E=\int\delta(H(\bm{q},\bm{p})-E)\dd{\bm{\Gamma}}, \quad \rho_E(\bm{q},\bm{p})=\frac{1}{\Sigma_E}\delta(H(\bm{q},\bm{p})-E)\]
可以看出在给定能量时，能量面上每个可及微观状态都被赋予相同权重。再考虑\textbf{混合流}，一种特殊的遍历流，\textcolor{blue}{\textbf{从任意时刻$t$的（非平衡）分布$\rho(\bm{q},\bm{p},t)$开始演化，经过足够长的时间后，系统的分布都会演化成微正则分布}}：
\[\lim_{t\to\infty}\rho(\bm{q},\bm{p},t)=\rho_E(\bm{q},\bm{p})\]
刘维尔定理（体积元守恒）与混合效应相结合表明：相空间中的无穷小体积元会经历无限的拉伸与折叠，并最终弥散分布于整个能量曲面上。因此\textcolor{blue}{\textbf{微观分布与初始状态无关}}，这就是经典统计力学的核心假设之一。\textcolor{blue}{\textbf{微观分布一旦确定，宏观量都可以视为相空间的期望值}}。

\subsection{微正则系综}
上一小节我们介绍了微正则系综，并给出了微正则配分函数$\Sigma_E$和微正则分布$\rho_E$的定义。显然我们可以知道在微正则系综中系统粒子数$N$和能量$E$是固定的。我们再来考虑体积$V$，如果体积$V$是可变的那么，会导致哈密顿量的势能显含时间，不满足微正则系综的条件，因此微正则系综的控制变量是$(N,V,E)$。因此微正则系综既不与外界交换热，也不交换粒子，所有满足给定能量壳层的微观状态被赋予等概率。记微正则配分函数为：
\[\Omega(N,V,E)=\int\delta(H(\bm{q},\bm{p})-E)\dd{\bm{\Gamma}}\]
$\Omega(N,V,E)$可以认为是体系的微观状态数，微正则系综中每个可及状态的概率都等于$1/\Omega$，Boltzmann 熵定义为：
\[S(N,V,E)=k_{\rm B}\ln \Omega(N,V,E)\]
上述 Boltzmann 熵可以认为是 Gibbs 熵在微正则系综中（等概率情况）的特殊情况，从 Gibbs 熵的定义出发：
\[S=-k_{\rm B}\int\rho(\bm{q},\bm{p})\ln\rho(\bm{q},\bm{p})\dd{\bm{\Gamma}}\]
其中密度分布函数$\rho(\bm{q},\bm{p})$可以认为在一个极小能量壳层$(\bm{q},\bm{p})\in \Gamma_{E}$内均匀分布：
\[\rho(\bm{q},\bm{p})=\frac{1}{\Omega(N,V,E)} \quad \Rightarrow \quad 
S=-k_{\rm B}\int_{\Gamma_{E}}\frac{1}{\Omega}\ln\frac{1}{\Omega}\dd{\bm{\Gamma}}=-k_{\rm B}\frac{1}{\Omega}\ln\frac{1}{\Omega}\int_{\Gamma_{E}}\dd{\bm{\Gamma}}=-k_{\rm B}\frac{1}{\Omega}\ln\frac{1}{\Omega}\cdot \Omega=k_{\rm B}\ln\Omega\]
同时我们也可以证明当 Gibbs 熵对分布函数做变分，在满足归一化约束的条件下达到极大值时，分布函数必然是微正则分布$\rho_E$，此时 Gibbs 熵就退化成 Boltzmann 熵。考虑对以下函数做变分：
\[\mathcal{S}=-k_{\rm B}\int\rho(\bm{q},\bm{p})\ln\rho(\bm{q},\bm{p})\dd{\bm{\Gamma}}-\lambda\left(\int\rho(\bm{q},\bm{p})\dd{\bm{\Gamma}}-1\right) \quad \Rightarrow \quad \delta\mathcal{S}=-k_{\rm B}\int\delta\rho(\bm{q},\bm{p})\left[\ln\rho(\bm{q},\bm{p})+1-\lambda\right]\dd{\bm{\Gamma}}=0\]
因此要求$\rho(\bm{q},\bm{p})$为常数，带入归一化条件：
\[\int \rho(\bm{q},\bm{p}) \dd{\bm{\Gamma}} = \rho(\bm{q},\bm{p})\int \dd{\bm{\Gamma}} = \rho(\bm{q},\bm{p})\Omega(N,V,E) = 1 \quad \Leftrightarrow \quad \rho(\bm{q},\bm{p}) = \frac{1}{\Omega(N,V,E)}\]
此时，在对$\mathcal{S}$做二阶变分：
\[\delta^2\mathcal{S}=-k_{\rm B}\int\frac{1}{2}\delta^2\rho(\bm{q},\bm{p})\cdot\frac{1}{\rho(\bm{q},\bm{p})}\dd{\bm{\Gamma}}<0\]
即😄$\mathcal{S}$在微正则分布处取得极大值。此外如果系统的密度分布满足$\rho=\rho_1\rho_2$，则\textcolor{blue}{\textbf{体系的熵$S$满足可加性}}：
\[S=-k_{\rm B}\int\rho\ln\rho\dd{\bm{\Gamma}}=-k_{\rm B}\left(\int\rho_1\ln\rho_1\dd{\bm{\Gamma}_1}\cdot\int\rho_2\dd{\bm{\Gamma}_2}+\int\rho_2\ln\rho_2\dd{\bm{\Gamma}_2}\cdot\int\rho_1\dd{\bm{\Gamma}_1}\right)=S_1+S_2\]

对 Boltzmann 熵做全微分：
\[\dd S=k_{\rm B}\left(\pdv{\ln\Omega}{N}\right)_{V,E}\dd N+k_{\rm B}\left(\pdv{\ln\Omega}{V}\right)_{N,E}\dd V+k_{\rm B}\left(\pdv{\ln\Omega}{E}\right)_{N,V}\dd E\]
另一方面，由热力学基本关系：
\[\dd U=T\dd S-P\dd V+\mu \dd N \quad \Rightarrow \quad \dd S=\frac{1}{T}\dd U+\frac{P}{T}\dd V-\frac{\mu}{T}\dd N\]
在微正则系综里$U=E$，逐项比较便得到：
\begin{equation}\label{eq:nve_terms}
\frac{1}{T}=k_{\rm B}\left(\pdv{\ln\Omega}{E}\right)_{N,V}, \quad P=k_{\rm B} T\left(\pdv{\ln\Omega}{V}\right)_{N,E}, \quad \mu=-k_{\rm B} T\left(\pdv{\ln\Omega}{N}\right)_{V,E}
\end{equation}
一旦$T,P,\mu$被确定，其余常见热力学势都可以通过 Legendre 变换得到：
\begin{equation}\label{eq:Legendre_4}
\boxed{U=E,\quad F=U-TS=E-k_{\rm B} T\ln\Omega, \quad H=U+PV=E+PV, \quad G=F+PV=U+PV-TS}
\end{equation}
对于可加的均匀体系，Euler 关系还给出$G=\sum_i\mu_i N_i$，其中$i$指代不同种类的粒子。
\begin{proposition}[微正则系综下的单原子理想气体]
在体积$V$的盒子中的$N$个单原子理想气体，体系哈密顿量满足：
\[H=\sum_{i=1}^{N}\frac{\bm{p}_i^2}{2m}=E\]
\end{proposition}
按照定义，容易看出微正则配分函数为：
\[\widetilde{\Omega}(N,V,E)=\frac{1}{h^{3N}}\int \delta\!\left(\sum_{i=1}^{N}\frac{\bm{p}_i^2}{2m}-E\right)\dd^{3N}\bm{q}\,\dd^{3N}\bm{p}=\frac{V^N}{h^{3N}}\int \delta\!\left(\frac{\bm{p}^2}{2m}-E\right)\dd^{3N}\bm{p}\equiv \frac{V^N}{h^{3N}}I\]
上式中的动量积分可以实际上是$3N$维空间中对一个半径为$\sqrt{2mE}$的超球面进行积分。利用$\delta$函数的性质：
\[\delta\!\left(\frac{p^2}{2m}-E\right)=\frac{\delta(p-\sqrt{2mE})}{|f'(\sqrt{2mE})|}=\sqrt{\frac{m}{2E}}\delta(p-\sqrt{2mE})\]
这个超球面的面积为：
\[\begin{aligned}
I&=\int\delta\!\left(\frac{p^2}{2m}-E\right)p^{3N-1}\dd p\dd{\bm{\Gamma}_{3N-1}}=\frac{2\pi^{3N/2}}{\Gamma(3N/2)}\int_0^\infty p^{3N-1}\delta\!\left(\frac{p^2}{2m}-E\right)\dd p\\
&=\frac{2\pi^{3N/2}}{\Gamma(3N/2)}\sqrt{\frac{m}{2E}}\int_0^\infty p^{3N-1}\delta(p-\sqrt{2mE})\dd p=\frac{2\pi^{3N/2}}{\Gamma(3N/2)}\sqrt{\frac{m}{2E}}\sqrt{2mE}^{3N-1}=\frac{(2\pi m)^{3N/2}}{\Gamma(3N/2)}E^{3N/2-1}
\end{aligned}\]
从而：
\[\widetilde{\Omega}(N,V,E)=\frac{V^N}{h^{3N}}\frac{(2\pi m)^{3N/2}}{\Gamma(3N/2)}E^{3N/2-1}\]
利用 Stirling 公式$\ln\Gamma(n)\sim n\ln n-n$，对应的微正则熵为：
\[\widetilde{S}(N,V,E)=k_{\rm B}\ln \widetilde{\Omega}(N,V,E)\sim Nk_{\rm B}\left\{\ln\left[V\left(\frac{4\pi mE}{3Nh^2}\right)^{3/2}\right]+\frac{3}{2}\right\}\]
进而可以求解其他热力学函数，首先是温度$T$和压强$P$：
\[\frac{1}{T}=\left(\pdv{\widetilde{S}}{E}\right)_{N,V}\sim\frac{3Nk_{\rm B}}{2E}, \quad P=T\left(\pdv{\widetilde{S}}{V}\right)_{N,E}=\frac{Nk_{\rm B} T}{V}\]
稍微改写一下就可以获得单原子理想气体能量$U=E=3Nk_{\rm B} T/2$以及状态方程$PV=Nk_{\rm B}T$。计算化学势：
\[\widetilde{\mu}=-T\left(\pdv{\widetilde{S}}{N}\right)_{V,E}=-k_{\rm B} T\ln\left[V\left(\frac{4\pi mE}{3Nh^2}\right)^{3/2}\right]\]
这个结果已经暴露出问题：\textcolor{blue}{\textbf{化学势本应是强度量，但这里它直接依赖于总体积$V$而不是粒子密度$N/V$}}。更严重的问题是熵不再是广延量：考虑两份完全相同的同种理想气体，各自处于$(N,V,E)$。抽去隔板前总熵为$2\widetilde{S}(N,V,E)$，抽去隔板后若仍把粒子当成可分辨粒子来数，则$\widetilde{S}_{\rm after}=2\widetilde{S}(N,V,E)+2Nk_{\rm B}\ln 2$，\textcolor{blue}{\textbf{但这显然不合理，因为两边本来就是同一种气体、相同温度、相同密度，抽去隔板并没有产生新的可观测宏观状态，熵增理应为零}}。

上述问题也被称为\textbf{ Gibbs 悖论}。在20世纪初期 Gibbs 本人提出了一个备受争议的解决方案：如果两个完全相同的相交换粒子不改变围观状态，因此在围观状态数中引入了因子$1/N!$以消除过度计数问题：
\begin{equation}\label{eq:partition_nve_mono}
\Omega(N,V,E)=\frac{\widetilde{\Omega}(N,V,E)}{N!}=\frac{V^N}{N!h^{3N}}\frac{(2\pi m)^{3N/2}}{\Gamma(3N/2)}E^{3N/2-1}
\end{equation}
1911 年到 1916 年，Sackur、Tetrode、Planck 更进一步，将上述 Gibbs 提出的修正过的围观状态数应用于熵与化学式熵，解决了熵无法广延 [Sackur-Tetrode 公式] 和化学势不是强度量的问题：
\[\begin{aligned}
S(N,V,E) &=Nk_{\rm B}\left\{\ln\left[\frac{V}{N}\left(\frac{4\pi mE}{3Nh^2}\right)^{3/2}\right]+\frac{5}{2}\right\}\\
\mu(N,V,E)&=-k_{\rm B} T\ln\left[\frac{V}{N}\left(\frac{4\pi mE}{3Nh^2}\right)^{3/2}\right]
\end{aligned}\]
但这在当时并没有被所有人接受。1921 年 Ehrenfest 和 Trkal 还批评说：在经典力学里粒子原则上可以靠轨迹区分，所以$1/N!$并不是由经典动力学“强迫”出来的，而更像是为了让统计热力学与经验热力学一致而加上的规则。

但不管怎么说我们最终获得了正确的微正则配分函数，由它再来计算与粒子数相关的热力学量，就会得到完全正确且具有正常标度性质的结果。我们把单原子理想气体能量$U=E=3Nk_{\rm B} T/2$带入上述熵和化学式中对数函数中的那一大坨：
\[\frac{4\pi mE}{3Nh^2}=\frac{4\pi m}{3h^2}\left(\frac{3Nk_{\rm B} T}{2}\right)=\frac{2\pi m k_{\rm B} T}{h^2} \quad \Rightarrow \quad \Lambda=\frac{h}{\sqrt{2\pi mk_{\rm B}T}}\]
通过上述方式定义出来的$\Lambda$被称为\textbf{热波长}，它是一个长度量纲的物理量，代表了粒子热运动的空间尺度。之后我们会经常见到它。\textcolor{blue}{\textbf{记住，在微正则系综可以作为变量的函数只有N，V和E，因此即使上述热波长引入了温度也不要把温度作为变量去运算，正确的做法应该是以N，V和E作为变量运算后再最后带入热波长进行化简！！！}}最后汇总一下微正则系综下单原子理想气体的热力学函数，记粒子数密度$n=N/V$：
\begin{theorem}[微正则系综下单原子理想气体的热力学函数]
\[\Omega(N,V,E)=\frac{V^N}{N!h^{3N}}\frac{(2\pi m)^{3N/2}}{\Gamma(3N/2)}E^{3N/2-1}, \quad S(N,V,E) \sim Nk_{\rm B}\left[\frac{5}{2}-\ln(n\Lambda^3)\right]\]
\[\left\{\begin{aligned}
U&=E=\frac{3}{2}Nk_{\rm B} T\\
F&=U-TS=-Nk_{\rm B} T\left[\ln\left(\frac{V}{N\Lambda^3}\right)+1\right]\\
H&=U+PV=\frac{5}{2}Nk_{\rm B} T\\
G&=F+PV=\mu N=Nk_{\rm B} T\ln(n\Lambda^3)\\
\end{aligned}\right.
\quad \Rightarrow \quad
\left\{\begin{aligned}
PV&=Nk_{\rm B} T\\
C_V&=\left(\pdv{U}{T}\right)_{V,N}=\frac{3}{2}Nk_{\rm B}\\
C_P&=\left(\pdv{H}{T}\right)_{P,N}=\frac{5}{2}Nk_{\rm B}\\
\mu&=k_{\rm B} T\ln(n\Lambda^3)
\end{aligned}\right.\]
\end{theorem}

最后的最后我们以马后炮角度来看看 Gibbs 悖论的解决方案。作为一个21世纪的学过量子力学的人我们知道$1/N!$的引入是由于\textcolor{blue}{\textbf{微观粒子的全同性}}的必然要求。Gibbs 在一个还没有量子力学的时代敢提出这个解决方案也是非常的超前。但我们也要注意到，量子力学的全同性原则是一个后续发展出来的理论，而不是从经典力学中直接推导出来的，因此在经典力学框架下，$1/N!$的引入确实是一个人为的修正。

但是话说回来，原先在经典统计中定义的相空间微元：
\[\dd{\bm{\Gamma}}=\frac{\dd{^{3N}\bm{p}}\dd{^{3N}\bm{q}}}{h^{3N}}\]
在经典框架内也并不意味着有问题。比如在很多教材中会出现的\textbf{经典定域子系统}，如果考虑经典粒子的话上述定义是完全合理的。但是当我们把统计的边界逐渐推向微观世界时，经典理论的框架就会逐渐失效。其实把相点的体积取为 Planck 常数也是为了让经典统计力学的结果与量子统计力学的结果在热力学极限下保持一致而引入的一个修正。这也是在经典框架下迫近微观量子世界时不得不做的妥协。

\newpage
\subsection{正则系综与玻尔兹曼分布}
从时间上来说，人类对统计力学的研究始于 Boltzmann 对微正则系综（孤立系统）中能量交换的研究，因此我们也将以 Boltzmann 分布的推导作开始本小节正则系综的介绍。考虑一个小系统$A$与一个巨大热浴$B$接触，二者组成的总体系孤立，总能量固定为$E_{\rm tot}$，是一个微正则系综。若系统$A$处于能量本征值$E_i$对应的微观状态，则热浴的能量就是$E_{\rm tot}-E_i$，于是该状态出现的概率正比于热浴在该能量下的微观态数：
\[P_i\propto \Omega_B(E_{\rm tot}-E_i)\]
给定体系能量后热浴$B$的能量也确定，因此也是微正则系综，对热浴态数在$E_{\rm tot}$附近做 Taylor 展开：
\[\ln \Omega_B(E_{\rm tot}-E_i)=\ln \Omega_B(E_{\rm tot})-\left(\pdv{\ln\Omega_B}{E_B}\right)_{V,N}E_i+\cdots \quad \Rightarrow \quad \left(\pdv{\ln\Omega_B}{E_B}\right)_{V,N}=\frac{1}{k_{\rm B}}\left(\pdv{S_B}{E_B}\right)_{V,N}=\frac{1}{k_{\rm B} T}\]
进而可知系统$A$处于能量$E_i$的概率（简并度为$g_i$）以及归一化因子$Z$为：
\[P_i=\frac{g_i}{Z}\exp\left(-\frac{E_i}{k_{\rm B} T}\right)\equiv\frac{g_i}{Z}e^{-\beta E_i}, \quad Z=\sum_i g_i\exp\left(-\frac{E_i}{k_{\rm B} T}\right)\equiv\sum_i g_i e^{-\beta E_i}\]
Boltzmann 研究的系统能与热浴交换能量、但体积$V$和粒子数$N$都固定，其全部微观状态按这一权重组成的统计集合就称为\textcolor{blue}{\textbf{正则系综}}。因此正则系综的自然变量是$(N,V,T)$。在经典相空间表述下，玻尔兹曼分布可以写成：
\[\rho(\xi)=\frac{e^{-\beta H(\xi)}}{Z(N,V,T)},\quad Z(N,V,T)=\int e^{-\beta H(\xi)}\dd\bm{\Gamma}\]
从 Gibbs 熵定义出发，带入正则分布可以获得熵关于正则配分函数的表达式：
\[S=-k_{\rm B}\int \rho(\xi)\ln\rho(\xi)\dd\Gamma=k_{\rm B}\beta\int \rho(\xi)H(\xi)\dd\Gamma+k_{\rm B}\ln Z(N,V,T)\int \rho(\xi)\dd\Gamma=k_{\rm B}\beta U+k_{\rm B}\ln Z(N,V,T)\]

从微正则系综的角度回看正则系综，还能得到一个很有启发的结果。若把总概率改写成能量的分布，则：
\begin{equation}\label{eq:eqnvt}
p(E)=\frac{\Omega(N,V,E)e^{-\beta E}}{Z(N,V,T)}=\frac{1}{Z(N,V,T)}\exp\left[\frac{S(N,V,E)}{k_{\rm B}}-\beta E\right]=\frac{1}{Z(N,V,T)}e^{-\beta F(N,V,E)}
\end{equation}
因此最概然能量对应的其实就是$F(N,V,E)$的极小值点，这也解释了为什么\textcolor{blue}{\textbf{ Helmholtz 自由能会成为正则系综的核心热力学势}}。由定义（callback 一下$\xi=(\bm{p},\bm{q})$是相点）出发推导其他热力学函数，如内能：
\[U=\int H(\xi)\rho(\xi)\dd\Gamma=\frac{1}{Z}\int H(\xi)e^{-\beta H(\xi)}\dd\Gamma=-\frac{1}{Z}\pdv{Z}{\beta}=-\pdv{\ln Z}{\beta}\]
再利用$F=U-TS$，得到 Helmholtz 自由能$F=U-T\left[k_{\rm B}\beta U+k_{\rm B}\ln Z\right]=-k_{\rm B}T\ln Z$。考虑$F$的全微分$\dd F=-S\dd T-P\dd V+\mu \dd N$可以得到压强$P$和化学势$\mu$，以及其他热力学函数如下：
\begin{theorem}[正则系综的基本热力学关系]
\[\rho(\xi)=\frac{e^{-\beta H(\xi)}}{Z},\quad Z=\int e^{-\beta H(\xi)}\dd\Gamma, \quad S=k_{\rm B}\left(\ln Z+\beta U\right)\]
\[\left\{\begin{aligned}
U&=-\pdv{\ln Z}{\beta}\\
F&=-k_{\rm B} T\ln Z\\
H&=-\pdv{\ln Z}{\beta}+k_{\rm B}TV\left(\pdv{\ln Z}{V}\right)_{N,T}\\
G&=-k_{\rm B} T\ln Z+k_{\rm B}TV\left(\pdv{\ln Z}{V}\right)_{N,T}
\end{aligned}\right.
\quad \Rightarrow \quad
\left\{\begin{aligned}
P&=k_{\rm B} T\left(\pdv{\ln Z}{V}\right)_{N,T}\\
\mu&=-k_{\rm B} T\left(\pdv{\ln Z}{N}\right)_{V,T}\\
C_V&=\left(\pdv{U}{T}\right)_{N,V}=k_{\rm B}\beta^2\pdv[2]{\ln Z}{\beta}\\
\expval{(\Delta E)^2}&=\pdv[2]{\ln Z}{\beta}=k_{\rm B} T^2 C_V
\end{aligned}\right.\]
\end{theorem}
如果体系满足通常的可加性关系，则依然有$G=\mu N$。正则系综里一个非常重要的结果是能量涨落公式：
\[\expval{(\Delta E)^2}=\expval{H^2}-\expval{H}^2=\frac{1}{Z}\pdv[2]{Z}{\beta}-\frac{1}{Z^2}\left(\pdv{Z}{\beta}\right)^2=\pdv{\beta}\left(\frac{1}{Z}\pdv{Z}{\beta}\right)=\pdv[2]{\ln Z}{\beta}=k_{\rm B} T^2 C_V\]
因此热容不仅描述“升温有多难”，也直接控制平衡能量涨落的大小。下面我们来看两个具体例子。

\begin{proposition}[正则系综下的单原子理想气体]
在在温度$T$，体积$V$的盒子中的$N$个单原子理想气体，体系哈密顿量由动能给出：
\[H=\sum_{i=1}^{N}\frac{\bm{p}_i^2}{2m}=\sum_{i=1}^{N}\frac{1}{2}m\bm{v}_i^2=\frac{1}{2}m\bm{v}^2, \quad \bm{v}^2=\sum_{i=1}^{N}\bm{v}_i^2\]
\end{proposition}
对单个理想气体粒子而言，玻尔兹曼权重直接给出三维速度分布，积分掉角度部分，另$v=|\bm{v}|$，则：
\[f(\bm{v})=\left(\frac{m}{2\pi k_{\rm B} T}\right)^{3/2}\exp\left(-\frac{m\bm{v}^2}{2k_{\rm B} T}\right) \quad \Rightarrow \quad f(v)=4\pi \left(\frac{m}{2\pi k_{\rm B} T}\right)^{3/2}v^2\exp\left(-\frac{mv^2}{2k_{\rm B} T}\right)\]
把变量从$v$换成$\varepsilon=mv^2/2$，由概率守恒$f(\varepsilon)\dd\varepsilon=f(v)\dd v$可得单粒子平动能量分布：
\[f(\varepsilon)
=f(v)\left|\dv{v}{\varepsilon}\right|=4\pi \left(\frac{m}{2\pi k_{\rm B} T}\right)^{3/2}\left(\frac{2\varepsilon}{m}\right)\exp\left(-\frac{\varepsilon}{k_{\rm B} T}\right)\frac{1}{\sqrt{2m\varepsilon}}=\frac{2}{\sqrt{\pi}} \frac{\sqrt{\varepsilon}}{(k_{\rm B} T)^{3/2}}e^{-\varepsilon/(k_{\rm B} T)}\]
最概然能量由$f(\varepsilon)$的极大值条件决定。对$\ln f(\varepsilon)$求导更方便：
\[
\dv{}{\varepsilon}\ln f(\varepsilon)=\frac{1}{2\varepsilon}-\frac{1}{k_{\rm B} T}=0
\quad \Rightarrow \quad
\varepsilon_{\rm mp}=\frac{1}{2}k_{\rm B} T.
\]
更一般地，令$x=\varepsilon/(k_{\rm B}T)$，则任意矩都可以写成 Gamma 函数积分：
\[\expval{\varepsilon^n}=\int_0^\infty \varepsilon^n f(\varepsilon)\dd\varepsilon=\frac{2}{\sqrt{\pi}}(k_{\rm B} T)^n\int_0^\infty x^{n+1/2}e^{-x}\dd x=\frac{2}{\sqrt{\pi}}\Gamma\left(n+\frac{3}{2}\right)(k_{\rm B} T)^n\]
取$n=1,2$便得到平均能量$\expval{\varepsilon}$和平均平方能量$\expval{\varepsilon^2}$，以及能量涨落$\expval{(\Delta \varepsilon)^2}$：
\[\boxed{\expval{\varepsilon}=\frac{3}{2}k_{\rm B} T,\quad \expval{\varepsilon^2}=\frac{15}{4}(k_{\rm B} T)^2,\quad \expval{(\Delta \varepsilon)^2}=\expval{\varepsilon^2}-\expval{\varepsilon}^2=\frac{3}{2}(k_{\rm B} T)^2}\]

上式给出的是\textbf{单粒子}能量$\varepsilon$的统计性质。由于理想气体粒子没有相互作用，所以在正则系综中$\varepsilon_1,\cdots,\varepsilon_N$彼此独立同分布，并且都服从上面的$f(\varepsilon)$。因此总能量的平均值可以直接由单粒子平均能量相加得到：
\[E=\sum_{i=1}^{N}\varepsilon_i \quad \Rightarrow \quad\expval{E}=\sum_{i=1}^{N}\expval{\varepsilon_i}=N\expval{\varepsilon}=\frac{3N}{2}k_{\rm B} T\]
同理，独立随机变量的方差也会直接相加，于是：
\[\expval{(\Delta E)^2}=\sum_{i=1}^{N}\expval{(\Delta \varepsilon_i)^2}=N\expval{(\Delta \varepsilon)^2}=\frac{3N}{2}(k_{\rm B} T)^2\]
如果还想求$\expval{E^2}$，就把：
\[\expval{E^2}=\left\langle\sum_{i=1}^{N}\varepsilon_i\right\rangle^2=\sum_{i=1}^{N}\left\langle\varepsilon_i\right\rangle^2+2\sum_{i<j}\left\langle\varepsilon_i\varepsilon_j\right\rangle=N\expval{\varepsilon^2}+N(N-1)\expval{\varepsilon}^2=\frac{3N}{2}\left(\frac{3N}{2}+1\right)(k_{\rm B} T)^2\]
\textcolor{blue}{\textbf{不过知道$\expval{\varepsilon}$、$\expval{\varepsilon^2}$这类单粒子矩，还不足以直接求出全部热力学函数}}，我们还是要从相空间理论出发，从单原子理想气体微正则系综配分函数 [Eq.\eqref{eq:partition_nve_mono}] 和正则系综能量分布 [Eq.\eqref{eq:eqnvt}]：
\[P(E)=\frac{\Omega(N,V,E)e^{-\beta E}}{Z(N,V,T)}, \quad \Omega(N,V,E)=\frac{V^N}{N!h^{3N}}\frac{(2\pi m)^{3N/2}}{\Gamma(3N/2)}E^{3N/2-1}\]
因此正则配分函数$Z(N,V,T)$为：
\[\begin{aligned}
Z(N,V,T)&=\int_0^\infty \Omega(N,V,E)e^{-\beta E}\dd E=\frac{V^N}{N!h^{3N}}\frac{(2\pi m)^{3N/2}}{\Gamma(3N/2)}\int_0^\infty E^{3N/2-1}e^{-\beta E}\dd E\\
&=\frac{V^N}{N!h^{3N}}\frac{(2\pi m)^{3N/2}}{\Gamma(3N/2)}\beta^{-3N/2}\int_0^\infty x^{3N/2-1}e^{-x}\dd x=\frac{V^N}{N!h^{3N}}\frac{(2\pi m)^{3N/2}}{\beta^{3N/2}}=\frac{1}{N!}\left(\frac{V}{\Lambda^3}\right)^N
\end{aligned}\]
把正则配分函数$Z(N,V,T)$代回正则能量分布$P(E)$，发现理想气体总动能便服从$\Gamma$分布：
\[P(E)=\frac{\beta^{3N/2}}{\Gamma(3N/2)}E^{3N/2-1}e^{-\beta E} \quad \Rightarrow \quad \expval{E^n}=\int_0^\infty E^n P(E)\dd E=\frac{\Gamma(3N/2+n)}{\Gamma(3N/2)}\frac{1}{\beta^n}\]
\[\boxed{\expval{E}=\frac{3N}{2\beta}=\frac{3N}{2}k_{\rm B} T,\quad \expval{E^2}=\frac{3N}{2}\left(\frac{3N}{2}+1\right)(k_{\rm B} T)^2, \quad \expval{(\Delta E)^2}=\expval{E^2}-\expval{E}^2=\frac{3N}{2}(k_{\rm B} T)^2}\]

\begin{theorem}[正则系综下单原子理想气体的热力学函数]
\[Z(N,V,T)=\frac{1}{N!}\left(\frac{V}{\Lambda^3}\right)^N,\quad S=k_{\rm B}\left(\ln Z+\beta U\right)\sim Nk_{\rm B}\left[\ln\left(\frac{V}{N\Lambda^3}\right)+\frac{5}{2}\right]\]
\[\left\{\begin{aligned}
F&=-k_{\rm B}T\ln Z\sim -Nk_{\rm B} T\left[\ln\left(\frac{V}{N\Lambda^3}\right)+1\right]\\
U&=-\pdv{\ln Z}{\beta}=\frac{3}{2}Nk_{\rm B} T\\
H&=U+PV=\frac{5}{2}Nk_{\rm B} T\\
G&=F+PV=\mu N=Nk_{\rm B} T\ln\left(\frac{N\Lambda^3}{V}\right)
\end{aligned}\right.
\quad \Rightarrow \quad
\left\{\begin{aligned}
P&=k_{\rm B} T\left(\pdv{\ln Z}{V}\right)_{N,T}=\frac{Nk_{\rm B} T}{V}\\
\mu&=\left(\pdv{F}{N}\right)_{V,T}=k_{\rm B} T\ln\left(\frac{N\Lambda^3}{V}\right)\\
C_V&=\left(\pdv{U}{T}\right)_{N,V}=\frac{3}{2}Nk_{\rm B}\\
C_P&=\left(\pdv{H}{T}\right)_{N,P}=\frac{5}{2}Nk_{\rm B}
\end{aligned}\right.\]
\end{theorem}

\begin{proposition}[正则系综下的两态系统]
设有$N$个彼此独立、固定在晶格位置上的缺陷或杂质，每个缺陷只有基态和激发态两种可能，记占据数$n_i=0,1$，激发能为$\varepsilon$，则哈密顿量为：
\[H=\varepsilon\sum_{i=1}^{N}n_i\]
由于这些缺陷是固定在各自位置上的，所以它们在经典计数里是可区分的。
\end{proposition}
此时配分函数可以直接因子化：
\[Z(N,V,T)=\sum_{\{n_i\}}\exp\left(-\beta\varepsilon\sum_i n_i\right)=\prod_{i=1}^{N}\left(\sum_{n_i=0}^{1}e^{-\beta\varepsilon n_i}\right)=\left(1+e^{-\beta\varepsilon}\right)^N\]
继而可以构建出正则系综下相空间的密度分布函数：
\[\rho(\{n_i\})=\frac{e^{-\beta H(\{n_i\})}}{Z(N,V,T)}=\frac{1}{Z(N,V,T)}\exp\left(-\beta\varepsilon\sum_{\{n_i\}} n_i\right)\]
在位置$\alpha$上粒子被被占据的概率为：
\[p_\alpha=\frac{1}{Z(N,V,T)}\sum_{\{n_i\}}n_\alpha\exp\left(-\beta\varepsilon\sum_i n_i\right)=\frac{e^{-\beta\varepsilon}}{1+e^{-\beta\varepsilon}}=\frac{1}{1+e^{\beta\varepsilon}}\]
\begin{theorem}[正则系综下两态系统的热力学函数]
\[Z(N,V,T)=\left(1+e^{-\beta\varepsilon}\right)^N\]
\[\left\{\begin{aligned}
F&=-k_{\rm B}T\ln Z=-Nk_{\rm B} T\ln\left(1+e^{-\beta\varepsilon}\right)\\
U&=-\pdv{\ln Z}{\beta}=\frac{N\varepsilon}{e^{\beta\varepsilon}+1}\\
H&=U+PV=U=\frac{N\varepsilon}{e^{\beta\varepsilon}+1}\\
G&=F+PV=F=-Nk_{\rm B} T\ln\left(1+e^{-\beta\varepsilon}\right)
\end{aligned}\right.
\quad \Rightarrow \quad
\left\{\begin{aligned}
\bar n_{\rm exc}&=Np_{\alpha}=\frac{N}{1+e^{\beta\varepsilon}}\\
P&=k_{\rm B} T\left(\pdv{\ln Z}{V}\right)_{N,T}=0\\
S&=Nk_{\rm B}\left[\ln\left(1+e^{-\beta\varepsilon}\right)+\frac{\beta\varepsilon}{e^{\beta\varepsilon}+1}\right]\\
C_V&=\left(\pdv{U}{T}\right)_{N,V}=Nk_{\rm B}(\beta\varepsilon)^2\frac{e^{\beta\varepsilon}}{(e^{\beta\varepsilon}+1)^2}
\end{aligned}\right.\]
\end{theorem}

在正则系综中通过对每一个给定能量的微正则配分函数$\Omega(N,V,E)$进行 Boltzmann 权重加权，我们得到了能量分布$p(E)$，进而构建了正则系综的配分函数$Z(N,V,T)$。巨正则系综的思路也是类似的，在正则系综的配分函数$Z(N,V,T)$的基础上对粒子数$N$进行求和。

\newpage
\subsection{巨正则系综}
在上一小节的基础上，考虑一个小系统$A$，其包含$N_i$个粒子且具有能量$E_i$，我们认为小系统$A$可以与热浴$B$交换能量与粒子，那小系统$A$出现状态$i$的概率将正比于此时热浴的微观状态数：
\[P_i\propto \Omega_B(E_{\rm tot}-E_i, N_{\rm tot}-N_i)\]
对热浴态数在$E_{\rm tot}$和$N_{\rm tot}$附近做 Taylor 展开并且只取一阶项并带入微正则系综的物理量表达式 [Eq.\eqref{eq:nve_terms}]：
\[\ln \Omega_B(E_{\rm tot}-E_i)\approx\ln \Omega_B(E_{\rm tot})-\left(\pdv{\ln\Omega_B}{E_B}\right)_{V,N}E_i-\left(\pdv{\ln\Omega_B}{N_B}\right)_{V,E}N_i=\ln \Omega_B(E_{\rm tot})-\frac{E_i}{k_{\rm B} T}+\frac{\mu N_i}{k_{\rm B} T}\]
进而可知系统$A$处于能量$E_i$和粒子数$N_i$的概率（简并度为$g_i$）以及归一化因子$\Xi$为：
\[P_i=\frac{g_i}{\Xi}\exp\left(-\frac{E_i}{k_{\rm B} T}+\frac{\mu N_i}{k_{\rm B} T}\right)\equiv\frac{g_i}{\Xi}e^{-\beta E_i+\beta\mu N_i}, \quad \Xi(T,V,\mu)=\sum_i g_i\exp\left(-\frac{E_i}{k_{\rm B} T}+\frac{\mu N_i}{k_{\rm B} T}\right)\equiv\sum_i g_i e^{-\beta E_i+\beta\mu N_i}\]
这里脚标$i$表示了系统带有$N_i$个粒子和$E_i$能量的状态，用传统相空间理论的相点形式重写上式：
\[\rho(\xi)=\frac{e^{-\beta [H_{N_{\xi}}(\xi)-\mu N_{\xi}]}}{\Xi(T,V,\mu)},\quad \Xi(T,V,\mu)=\int e^{-\beta [H_{N_{\xi}}(\xi)-\mu N_{\xi}]}\dd\bm{\Gamma}\]
带入 Gibbs 熵定义：
\[S=-k_{\rm B}\int \rho\ln\rho\dd\Gamma=k_{\rm B}\left(\beta\int \rho(\xi)H_{N_{\xi}}(\xi)\dd\Gamma-\beta\mu\int \rho(\xi)N_{\xi}\dd\Gamma+\ln \Xi\int \rho(\xi)\dd\Gamma\right)=k_{\rm B}\left(\beta U-\beta\mu \expval{N}+\ln \Xi\right)\]

将相点元对粒子数进行分类$\dd{\bm{\Gamma}}=\dd{\bm{\Gamma}_E}\dd{\bm{\Gamma}_N}$，其中$\dd{\bm{\Gamma}_N}$粒子数空间的积分是在离散的计数测度空间上进行的，因此可以把巨正则分布函数$\rho(\xi)$看成是一个能量分布和粒子数分布的乘积：
\begin{equation}\label{eq:partition_nvmu}
\begin{aligned}
\Xi(T,V,\mu)&=\int e^{-\beta [H_{N_{\xi}}(\xi)-\mu N_{\xi}]}\dd\bm{\Gamma}_E\dd\bm{\Gamma}_N=\int\dd\bm{\Gamma}_N e^{\beta\mu N_{\xi}}\int e^{-\beta H_{N_{\xi}}(\xi)}\dd\bm{\Gamma}_E\\
&=\sum_{N=0}^{\infty}e^{\beta\mu N}\int e^{-\beta H_N(\xi)}\dd\bm{\Gamma}_E=\sum_{N=0}^{\infty}Z(N,V,T)e^{\beta\mu N}
\end{aligned}
\end{equation}
因此，在正则系综的基础上考虑不同粒子数，并且引入了新的控制变量化学势$\mu$，就得到了巨正则系综，其自然变量是$(T,V,\mu)$。

类似正则系综中的 Helmholtz 自由能，在巨正则系综中定义\textbf{巨势}$\Phi_{\rm G}(T,V,\mu)=-k_{\rm B}T\ln \Xi$来简化热力学函数的表达。由全微分$\dd\Phi_{\rm G}=-S\dd T-P\dd V-\expval{N}\dd\mu$可以得到平均粒子数$\expval{N}$、压强$P$与熵$S$（也可从定义获得）。其他热力学函数可以通过 Legendre 变换得到 [Eq.\eqref{eq:Legendre_4}]：
\begin{theorem}[巨正则系综的基本热力学关系]
\[\rho(\xi)=\frac{e^{-\beta [H_{N_{\xi}}(\xi)-\mu N_{\xi}]}}{\Xi(T,V,\mu)},\quad \Xi(T,V,\mu)=\int e^{-\beta [H_{N_{\xi}}(\xi)-\mu N_{\xi}]}\dd\bm{\Gamma}\]
\[\left\{\begin{aligned}
\Phi_{\rm G}&=-k_{\rm B}T\ln\Xi\\
U&=-\left(\pdv{\ln\Xi}{\beta}\right)_{V,\mu}+\mu \expval{N}\\
F&=\Phi_{\rm G}+\mu \expval{N}=-k_{\rm B}T\ln\Xi+\mu \expval{N}\\
H&=U+PV\\
G&=F+PV=\mu \expval{N}
\end{aligned}\right.
\quad \Rightarrow \quad
\left\{\begin{aligned}
\expval{N}&=-\left(\pdv{\Phi_{\rm G}}{\mu}\right)_{T,V}=k_{\rm B}T\left(\pdv{\ln\Xi}{\mu}\right)_{T,V}\\
P&=-\left(\pdv{\Phi_{\rm G}}{V}\right)_{T,\mu}=k_{\rm B}T\left(\pdv{\ln\Xi}{V}\right)_{T,\mu}\\
S&=k_{\rm B}\left(\ln\Xi+\beta U-\beta\mu \expval{N}\right)\\
\expval{(\Delta N)^2}&=k_{\rm B}T\left(\pdv{\expval{N}}{\mu}\right)_{T,V}
\end{aligned}\right.\]
\end{theorem}
如果体系满足通常的可加性关系，则仍然有$G=\mu \expval{N}$。同时从熵出发，可得$TS=-\Phi_{\rm G}+U-\mu \expval{N}$，稍作变形可得$\Phi_{\rm G}=-PV$。上述公式中的粒子涨落$\expval{(\Delta N)^2}$可以另$\alpha=\beta\mu=\mu/k_{\rm B}T$，稍作运算：
\[\expval{(\Delta N)^2}=\expval{N^2}-\expval{N}^2=\frac{1}{\Xi}\pdv[2]{\Xi}{\alpha}-\frac{1}{\Xi^2}\left(\pdv{\Xi}{\alpha}\right)^2=\pdv{\alpha}\left(\frac{1}{\Xi}\pdv{\Xi}{\alpha}\right)=\pdv[2]{\ln\Xi}{\alpha}=(k_{\rm B}T)^2\pdv[2]{\ln\Xi}{\mu}=k_{\rm B}T\left(\pdv{\expval{N}}{\mu}\right)_{T,V}\]

\begin{proposition}[巨正则系综下的单原子理想气体]
在温度$T$下，体积为$V$的单原子理想气体具有化学势$\mu$，当体系中含有$N$个粒子时，体系能量为：
\[H=\sum_{i=1}^{N}\frac{\bm{p}_i^2}{2m}=\sum_{i=1}^{N}\frac{1}{2}m\bm{v}_i^2=\frac{1}{2}m\bm{v}^2, \quad \bm{v}^2=\sum_{i=1}^{N}\bm{v}_i^2\]
\end{proposition}
将单原子理想气体的正则配分函数：
\[Z(N,V,T)=\frac{1}{N!}\left(\frac{V}{\Lambda^3}\right)^N\]
代入巨正则配分函数的定义式 [Eq.\eqref{eq:partition_nvmu}]，可以得到巨正则配分函数$\Xi$的具体表达式：
\[\Xi(T,V,\mu)=\sum_{N=0}^{\infty}Z(N,V,T)e^{\beta\mu N}=\sum_{N=0}^{\infty}\frac{1}{N!}\left(\frac{V}{\Lambda^3}\right)^N e^{\beta\mu N}=\sum_{N=0}^{\infty}\frac{1}{N!}\left(\frac{Ve^{\beta\mu}}{\Lambda^3}\right)^N=\exp\left(\frac{Ve^{\beta\mu}}{\Lambda^3}\right)\]
因此体系包含$N$个粒子的概率为：
\[p(N)=\frac{Z(N,V,T)}{\Xi}e^{\beta\mu N}=\exp\left(-\frac{Ve^{\beta\mu}}{\Lambda^3}\right)\frac{1}{N!}\left(\frac{Ve^{\beta\mu}}{\Lambda^3}\right)^N\]
上述概率分布正是参数$Ve^{\beta\mu}/\Lambda^3$的 Poisson 分布。计算平均粒子数$\expval{N}$可以发现 Poisson 分布的参数就是$\expval{N}$：
\[\expval{N}=\sum_{N=0}^{\infty}Np(N)=\sum_{N=1}^{\infty}Np(N)=\exp\left(-\frac{Ve^{\beta\mu}}{\Lambda^3}\right)\sum_{N=1}^{\infty}\frac{1}{(N-1)!}\left(\frac{Ve^{\beta\mu}}{\Lambda^3}\right)^N=\frac{Ve^{\beta\mu}}{\Lambda^3}\]
\[\begin{aligned}
\expval{(\Delta N)^2}&=\sum_{N=0}^{\infty}(N-\expval{N})^2p(N)=\sum_{N=0}^{\infty}N^2p(N)-\expval{N}^2=\exp\left(-\frac{Ve^{\beta\mu}}{\Lambda^3}\right)\sum_{N=0}^{\infty}\frac{N^2}{N!}\left(\frac{Ve^{\beta\mu}}{\Lambda^3}\right)^N-\expval{N}^2\\
&=\exp\left(-\frac{Ve^{\beta\mu}}{\Lambda^3}\right)\sum_{N=1}^{\infty}\frac{1}{(N-1)!}\left(\frac{Ve^{\beta\mu}}{\Lambda^3}\right)^N+\exp\left(-\frac{Ve^{\beta\mu}}{\Lambda^3}\right)\sum_{N=2}^{\infty}\frac{1}{(N-2)!}\left(\frac{Ve^{\beta\mu}}{\Lambda^3}\right)^N-\expval{N}^2\\
&=\frac{Ve^{\beta\mu}}{\Lambda^3}+\left(\frac{Ve^{\beta\mu}}{\Lambda^3}\right)^2-\expval{N}^2=\frac{Ve^{\beta\mu}}{\Lambda^3}=\expval{N}
\end{aligned}\]
说实话上面两个推导只是复习了一下 Poisson 分布的性质，平均值和方差都等于分布的参数。

我们再来看看理想气体状态方程，考虑压强$P$：
\[P=k_{\rm B}T\left(\pdv{\ln\Xi}{V}\right)_{T,\mu}=k_{\rm B}T\left[\pdv{V}\left(\frac{Ve^{\beta\mu}}{\Lambda^3}\right)\right]_{T,\mu}=k_{\rm B}T\frac{e^{\beta\mu}}{\Lambda^3}=\frac{\expval{N}k_{\rm B}T}{V}\]
记平均粒子数密度$\bar n=\bar N/V=e^{\beta\mu}/\Lambda^3$，则其他巨正则系综下的单原子理想气体热力学函数如下给出：
\begin{theorem}[巨正则系综下单原子理想气体的热力学函数]
\[\Xi(T,V,\mu)=\exp\left(\frac{zV}{\Lambda^3}\right), \quad S=k_{\rm B}\left(\ln\Xi+\beta U-\beta\mu \expval{N}\right)=\expval{N} k_{\rm B}\left[\ln\left(\frac{V}{\expval{N}\Lambda^3}\right)+\frac{5}{2}\right]\]
\[\left\{\begin{aligned}
\Phi_G&=-k_{\rm B}T\ln\Xi=-\expval{N} k_{\rm B}T\\
U&=-\left(\pdv{\ln\Xi}{\beta}\right)_{V,\mu}+\mu \expval{N}=\frac{3}{2}\expval{N} k_{\rm B}T\\
H&=U+PV=\frac{5}{2}\expval{N} k_{\rm B}T\\
F&=\Phi_G+\mu \expval{N}=-\expval{N} k_{\rm B}T\left[\ln\left(\frac{V}{\expval{N}\Lambda^3}\right)+1\right]\\
G&=F+PV=\mu \expval{N}=\expval{N} k_{\rm B}T\ln(\bar n\Lambda^3)
\end{aligned}\right.
\quad \Rightarrow \quad
\left\{\begin{aligned}
&P=k_{\rm B}T\left(\pdv{\ln\Xi}{V}\right)_{T,\mu}=\frac{\expval{N} k_{\rm B}T}{V}\\
&\mu=k_{\rm B}T\ln(\bar n\Lambda^3)\\
&C_V=\left(\pdv{U}{T}\right)_{V,\expval{N}}=\frac{3}{2}\expval{N} k_{\rm B}\\
&C_P=\left(\pdv{H}{T}\right)_{P,\expval{N}}=\frac{5}{2}\expval{N} k_{\rm B}\\
&\expval{(\Delta N)^2}=\expval{N}
\end{aligned}\right.\]
\end{theorem}

可以看到，无论在微正则、正则还是巨正则系综里，单原子理想气体最终都回到同一个状态方程；区别只在于中间计算时最方便的控制变量分别是$(N,V,E)$、$(N,V,T)$和$(T,V,\mu)$。

\newpage
\section{量子统计力学}\label{sec:quantum_sm}
\subsection{相空间理论量子化}
在章节\ref{sec:qm_basic}中我们简单介绍了量子力学，这里不再赘述。应该要注意的，量子力学中我们并没有涉及到温度这个概念，我们只关注了体系可能存在的量子态以及各种算符的本征值和期望值，但是我们没有讨论过系统如何在这些量子态上分布，这也是量子统计需要解决的问题，因此可以认为量子力学是量子统计在$T=0 \ \rm{K}$的特殊情况。从量子力学中的算符 [Eq.\eqref{eq:operator_qm}] 出发，假设$\{\ket{\phi_{\alpha}}\}$是一组完备基，则算符$\hat{O}$可以展开为：
\[\hat{O}=\mathbbm{1}\hat{O}\mathbbm{1}=\sum_{\alpha\alpha'}\ket{\phi_{\alpha}}w(\alpha)\bra{\phi_{\alpha}}\hat{O}\ket{\phi_{\alpha'}}w(\alpha')\bra{\phi_{\alpha'}}\]
计算算符$\hat{O}$在态$\ket{\psi}$上的期望值$\bra{\psi}\hat{O}\ket{\psi}$：
\[\begin{aligned}
\langle\hat{O}\rangle=\bra{\psi}\hat{O}\ket{\psi}&=\sum_{\alpha\alpha'}\bra{\psi}\ket{\phi_{\alpha}}w(\alpha)\bra{\phi_{\alpha}}\hat{O}\ket{\phi_{\alpha'}}w(\alpha')\bra{\phi_{\alpha'}}\ket{\psi}=\sum_{\alpha\alpha'}w(\alpha)\bra{\phi_{\alpha}}\hat{O}\ket{\phi_{\alpha'}}w(\alpha')\bra{\phi_{\alpha'}}\ket{\psi}\bra{\psi}\ket{\phi_{\alpha}}\\
&=\sum_{\alpha}w(\alpha)\bra{\phi_{\alpha}}\hat{O}\left(\sum_{\alpha'}\ket{\phi_{\alpha'}}w(\alpha')\bra{\phi_{\alpha'}}\right)\ket{\psi}\bra{\psi}\ket{\phi_{\alpha}}=\sum_{\alpha}w(\alpha)\bra{\phi_{\alpha}}\hat{O}\ket{\psi}\bra{\psi}\ket{\phi_{\alpha}}\\
\end{aligned}\]
定义密度算符$\hat{\rho}$为：
\[\ket{\psi}=\sum_{\alpha}c_{\alpha}\ket{\phi_{\alpha}} \quad \Rightarrow \quad \boxed{\hat{\rho}=\ket{\psi}\bra{\psi}=\sum_{\alpha\alpha'}c_{\alpha}c^{\dagger}_{\alpha'}\ket{\phi_{\alpha}}\bra{\phi_{\alpha'}}:=\sum_{\alpha\alpha'}\rho_{\alpha\alpha'}\ket{\phi_{\alpha}}\bra{\phi_{\alpha'}}}\]
上式中可以称矩阵元为$\rho_{\alpha\alpha'}=c_{\alpha}c^{\dagger}_{\alpha'}$的矩阵$\bm{\rho}$为密度矩阵。由此在基底$\{\ket{\phi_{\alpha}}\}$下，期望值$\langle\hat{O}\rangle$可以写成：
\[\boxed{\langle\hat{O}\rangle=\sum_{\alpha}w(\alpha)\bra{\phi_{\alpha}}\hat{O}\ket{\psi}\bra{\psi}\ket{\phi_{\alpha}}=\sum_{\alpha}w(\alpha)\bra{\phi_{\alpha}}\hat{O}\hat{\rho}\ket{\phi_{\alpha}}=\Tr(\hat{O}\hat{\rho})=\Tr(\hat{\rho}\hat{O})}\]
密度算符总满足$\hat{\rho}^{\dagger}=\hat{\rho}$和$\Tr\hat{\rho}=1$，且$\forall \ket{\phi}\in\mathcal{H}$有$\bra{\phi}\hat{\rho}\ket{\phi}\ge 0$，即是半正定的厄米算符。如果$\ket{\psi}$是一个纯态，则$\Tr\hat{\rho}^2=\Tr\hat{\rho}$；如果$\ket{\psi}$是一个混态，则$\Tr\hat{\rho}^2<\Tr\hat{\rho}$。我们可以把 Gibbs 熵在量子统计中推广为：
\[\boxed{S=-k_{\rm B}\Tr(\hat{\rho}\ln\hat{\rho})}\]
出于方便我们要求$\mathcal{H}=\text{span}\{\ket{\phi_{\alpha}}\}$是正交归一的完备基，或者说某个算符的本征态（纯态），那么$\hat{\rho}$可以写成：
\[\hat{\rho}=\sum_{\alpha}p_{\alpha}\ket{\phi_{\alpha}}\bra{\phi_{\alpha}}, \quad p_{\alpha}\ge 0, \quad \sum_{\alpha}p_{\alpha}=1\]
则 Gibbs 熵可以写成：
\[S=-k_{\rm B}\Tr(\hat{\rho}\ln\hat{\rho})=-k_{\rm B}\sum_{\alpha}p_{\alpha}\ln p_{\alpha}\]
纯态时某个$p_{\alpha}=1$而其余全为$0$，因此$S=0$；完全无信息时$\hat{\rho}=\mathbbm{1}/\dim\mathcal{H}$，熵取最大值$S=k_{\rm B}\ln\dim\mathcal{H}$。
\begin{proposition}[单电子自旋系统]
在章节\ref{sec:spin}中我们已经介绍了三个 Pauli 矩阵$\sigma_x$、$\sigma_y$和$\sigma_z$是电子自旋空间的基底，因此我们可以把它们作为完备基来展开密度算符：
\[\hat{\rho}=\frac{1}{2}\left(\mathbbm{1}+\mathbf{P}\cdot\hat{\bm{\sigma}}\right), \quad \mathbf{P}=(P_x,P_y,P_z)\in\mathbb{R}^3, \quad |\mathbf{P}|\le 1\]
其中$\mathbf{P}$是 Bloch 向量。当$|\mathbf{P}|=1$时，$\hat{\rho}$是一个纯态；当$|\mathbf{P}|<1$时，$\hat{\rho}$是一个混态。布洛赫向量的方向表示自旋的平均取向，而模长则表示自旋的纯度。

一个自然而然的问题是：电子密度矩阵的自由度是多少？理论上四个矩阵元是复数应该有8个自由度，但是由于厄米性，非对角线满足共轭关系且对角线上是实数，只有四个自由度；然后由$\Tr\hat{\rho}=1$，最终只有3个自由度，且是实数，与 Bloch 向量对应。我们可以计算电子自旋的期望值$\langle\hat{\mathbf{S}}\rangle$：
\[\langle\hat{\mathbf{S}}\rangle=\frac{\hbar}{2}\Tr(\hat{\rho}\hat{\bm{\sigma}})=\frac{\hbar}{4}\Tr\left(\hat{\bm{\sigma}}+\mathbf{P}\cdot\hat{\bm{\sigma}}^2\right)=\frac{\hbar}{4}\left(0+2\mathbf{P}\right)=\frac{\hbar}{2}\mathbf{P}\]
callback 一下 Pauli 矩阵的性质$\hat{\sigma}^2_i=\mathbf{I}$，$\Tr\hat{\sigma}_i=0$。如果不知道也可以直接按照定义展开$2 \times 2$矩阵计算。
\end{proposition}

\begin{proposition}[其他基底下的单电子自旋系统]
上述例子中我们采用了$\sigma_z$的本征态作为密度算符的基底，也即$\ket{\uparrow}$和$\ket{\downarrow}$，其方向刚好与$z$轴对齐。如果我们考虑旋转坐标轴到以单位向量$\mathbf{n}=(\sin\theta\cos\phi,\sin\theta\sin\phi,\cos\theta)$为$z$轴，那么$\hat{\rho}$的基底就变成了$\hat{\mathbf{n}}\cdot\hat{\bm{\sigma}}$的本征态$\ket{\uparrow_{\mathbf{n}}}$和$\ket{\downarrow_{\mathbf{n}}}$，其与原先基底的关系是：
\[\ket{\uparrow_{\mathbf{n}}}=\cos\frac{\theta}{2}\ket{\uparrow}+e^{i\phi}\sin\frac{\theta}{2}\ket{\downarrow}, \quad \ket{\downarrow_{\mathbf{n}}}=-e^{-i\phi}\sin\frac{\theta}{2}\ket{\uparrow}+\cos\frac{\theta}{2}\ket{\downarrow}\]
其中$\theta$和$\phi$是单位向量$\mathbf{n}$在球坐标系中的极角和方位角。此时沿着$\mathbf{n}$正方向的投影算符的展开式为：
\[\hat{P}_{\uparrow_{\mathbf{n}}}=\ket{\uparrow_{\mathbf{n}}}\bra{\uparrow_{\mathbf{n}}}=\frac{1}{2}\left(\mathbbm{1}+\mathbf{n}\cdot\hat{\bm{\sigma}}\right)\]
因此沿着$\mathbf{n}$方向的自旋期望值为：
\[\langle\hat{P}_{\uparrow_{\mathbf{n}}}\rangle=\Tr(\hat{P}_{\uparrow_{\mathbf{n}}}\hat{\rho})=\frac{1}{4}\Tr[\left(\mathbbm{1}+\mathbf{n}\cdot\hat{\bm{\sigma}}\right)\left(\mathbbm{1}+\mathbf{P}\cdot\hat{\bm{\sigma}}\right)]=\frac{1}{4}\left(2+0+0+2\mathbf{n}\cdot\mathbf{P}\right)=\frac{1}{2}(1+\mathbf{n}\cdot\mathbf{P})\]
当然还是可以暴力展开$2 \times 2$矩阵计算。
\end{proposition}

如果系统可以拆分成两个子系统$A$和$B$，即总希尔伯特空间是子空间张量积$\mathcal{H}_{AB}=\mathcal{H}_A\otimes\mathcal{H}_B$，记$\{\ket{i^{(A)}}\}$和$\{\ket{\mu^{(B)}}\}$分别是$\mathcal{H}_A$与$\mathcal{H}_B$的正交归一完备基，则$\{\ket{i^{(A)}}\otimes\ket{\mu^{(B)}}\}$构成总空间$\mathcal{H}_{AB}$的正交归一完备乘积基底。对于彼此独立的复合系统，其总密度算符满足$\hat{\rho}_{AB}=\hat{\rho}_A\otimes\hat{\rho}_B$：
\[\hat{\rho}_A=\sum_i p_i\ket{i^{(A)}}\bra{i^{(A)}}, \quad \hat{\rho}_B=\sum_{\mu}q_{\mu}\ket{\mu^{(B)}}\bra{\mu^{(B)}} \quad \Rightarrow \quad \hat{\rho}_{AB}=\sum_{i,\mu}p_iq_{\mu}\left(\ket{i^{(A)}}\bra{i^{(A)}}\right)\otimes\left(\ket{\mu^{(B)}}\bra{\mu^{(B)}}\right)\]
可以验证此时系统的熵满足可加性：
\[S_{AB}=-k_{\rm B}\sum_{i,\mu}p_iq_{\mu}\ln(p_iq_{\mu})=-k_{\rm B}\sum_ip_i\ln p_i-k_{\rm B}\sum_{\mu}q_{\mu}\ln q_{\mu}=S_A+S_B\]
如果我们只关心其中一个子系统，而忽略另一个子系统的自由度，就需要对后者作偏迹并定义约化密度算符：
\[\boxed{\hat{\rho}_A=\Tr_B\hat{\rho}_{AB}=\sum_{\mu}\bra{\mu^{(B)}}\hat{\rho}_{AB}\ket{\mu^{(B)}}}, \qquad \boxed{\hat{\rho}_B=\Tr_A\hat{\rho}_{AB}=\sum_i\bra{i^{(A)}}\hat{\rho}_{AB}\ket{i^{(A)}}}\]
偏迹的结果与在子空间所选基底无关，并且满足：
\[\Tr_A(\hat{\rho}_A\hat{O}_A)=\Tr_{AB}\left[\hat{\rho}_{AB}\left(\hat{O}_A\otimes\mathbbm{1}_B\right)\right]\]
因此$\hat{\rho}_A$完整编码了子系统$A$上一切局域可观测量的统计信息。

上面的内容如果反过来说，对于任意纯态$\ket{\Psi^{(AB)}}\in\mathcal{H}_A\otimes\mathcal{H}_B$，总可以找到两个子空间$\mathcal{H}_A$和$\mathcal{H}_B$的两组正交归一基底$\{\ket{u_{\alpha}^{(A)}}\}$和$\{\ket{v_{\alpha}^{(B)}}\}$，使它写成施密特分解，满足$\lambda_{\alpha}\ge 0$，$\sum\lambda_{\alpha}=1$：
\[\ket{\Psi^{(AB)}}=\sum_{\alpha}\sqrt{\lambda_{\alpha}}\ket{u_{\alpha}^{(A)}}\otimes\ket{v_{\alpha}^{(B)}}, \quad \hat{\rho}_A=\sum_{\alpha}\lambda_{\alpha}\ket{u_{\alpha}^{(A)}}\bra{u_{\alpha}^{(A)}}, \quad \hat{\rho}_B=\sum_{\alpha}\lambda_{\alpha}\ket{v_{\alpha}^{(B)}}\bra{v_{\alpha}^{(B)}}\]
两个子系统拥有完全相同的非零本征值，熵也相等。

\begin{proposition}[两态系统与单模光场的耦合]
考虑一个两态系统$A$，基底为$\{\ket{g},\ket{e}\}$，以及一个只保留$0,1$光子子空间的单模光场$F$，基底为$\{\ket{0},\ket{1}\}$。若总系统处于纯态$\ket{\Psi}=\sqrt{p}\ket{e}\otimes\ket{0}+e^{i\phi}\sqrt{q}\ket{g}\otimes\ket{1}$，其中$p+q=1$，总密度算符为：
\[\hat{\rho}=p\ket{e}\otimes\ket{0}\bra{e}\otimes\bra{0}+q\ket{g}\otimes\ket{1}\bra{g}\otimes\bra{1}+e^{i\phi}\sqrt{pq}\ket{e}\otimes\ket{0}\bra{g}\otimes\bra{1}+e^{-i\phi}\sqrt{pq}\ket{g}\otimes\ket{1}\bra{e}\otimes\bra{0}\]
对光与物质场作偏迹可得两态系统的约化密度算符：
\[\hat{\rho}_A=\Tr_F\hat{\rho}=\sum_{i=0,1}\bra{i}\hat{\rho}\ket{i}=p\ket{e}\bra{e}+(1-p)\ket{g}\bra{g}\]
\[\hat{\rho}_F=\Tr_A\hat{\rho}=\sum_{i=g,e}\bra{i}\hat{\rho}\ket{i}=p\ket{0}\bra{0}+(1-p)\ket{1}\bra{1}\]
因此两子系统具有相同的约化熵$S_A=S_F=-k_{\rm B}\left[p\ln p+(1-p)\ln(1-p)\right]$。特别地，当$p=1/2$时，耦合态纠缠最强，熵达到最大值$S_A=S_F=k_{\rm B}\ln 2$；当$p=0$或$p=1$时，总态退化为乘积态，约化熵为$0$。
\end{proposition}

\newpage
\subsection{量子平衡分布与三个量子系综}
在经典统计中，相空间概率密度$\rho_{\rm cl}(\bm{q},\bm{p},t)$满足 Liouville 方程：
\[\pdv{\rho_{\rm cl}}{t}+\{\rho_{\rm cl},H\}=0 \quad \Leftrightarrow \quad \pdv{\rho_{\rm cl}}{t}=\{H,\rho_{\rm cl}\}\]
而在量子统计中，分布对象从相空间概率密度变成了希尔伯特空间上的密度算符$\hat{\rho}(t)$。因此我们首先要把经典 Liouville 方程量子化。若孤立体系处于纯态$\ket{\psi(t)}$，则它满足薛定谔方程：
\[i\hbar \pdv{}{t}\ket{\psi(t)}=\hat{H}\ket{\psi(t)}\]
定义纯态对应的密度算符$\hat{\rho}(t)=\ket{\psi(t)}\bra{\psi(t)}$：
\[\boxed{i\hbar \pdv{\hat{\rho}}{t}=i\hbar \left(\pdv{}{t}\ket{\psi}\right)\bra{\psi}+i\hbar \ket{\psi}\left(\pdv{}{t}\bra{\psi}\right)=\hat{H}\ket{\psi}\bra{\psi}-\ket{\psi}\bra{\psi}\hat{H}=\hat{H}\hat{\rho}-\hat{\rho}\hat{H}=[\hat{H},\hat{\rho}]}\]
这就是 \textbf{Liouville-von Neumann 方程}。若$\hat{H}$不显含时间，则其形式解为：
\[\hat{\rho}(t)=e^{-i\hat{H}t/\hbar}\hat{\rho}(0)e^{i\hat{H}t/\hbar}\]
可以看到\textcolor{blue}{\textbf{经典中的 Poisson 括号$\{A,B\}$在量子化后对应于交换子$[\hat{A},\hat{B}]/(i\hbar)$}}。在经典力学中，Liouville 定理告诉我们相空间流是不可压缩的，概率密度沿相轨线保持不变；在量子力学中，闭系统演化是幺正的，因此$\Tr\hat{\rho}=1$、$\hat{\rho}\ge 0$以及$\hat{\rho}$的本征值谱都保持不变，从而 von Neumann 熵$S=-k_{\rm B}\Tr(\hat{\rho}\ln\hat{\rho})$在精确闭系统演化下也是守恒的。

统计力学真正关心的不是任意时刻的非平衡幺正演化，而是在给定宏观约束下长期稳定的平衡分布。因此我们只研究\textbf{平衡态}（非平衡态太难了），即密度算符不再显式随时间变化的情形：
\[\pdv{\hat{\rho}}{t}=0 \quad \Rightarrow \quad[\hat{H},\hat{\rho}_{\rm eq}]=0, \quad \hat{H}\ket{n}=E_n\ket{n}\]
这说明平衡态密度算符与哈密顿量对易，即$\hat{\rho}=f(\hat{H})$。这正是经典平衡条件$\{H,\rho_{\rm cl}^{\rm eq}\}=0$的量子对应。于是经典三大系综在量子理论中的量子化，只不过是把经典相空间分布换成满足上述平衡条件的密度算符。

\begin{proposition}[量子化的微正则系综]
首先看微正则系综，自然变量是$(N,V,E)$。若体系能量被限制在$E$上，则微正则分布函数可以表示为：
\[\hat{\rho}_{\rm micro}(E)=\frac{\delta(E-\hat{H})}{\Omega}=\frac{1}{\Omega}\sum_n\delta(E-E_n)\ket{n}\bra{n}, \quad \Omega=\Tr\left[\delta(\hat{H}-E)\right]=\sum_n\delta(E-E_n)\]
可以验证，在上述定义的微正则配分函数满足等概率原理：
\[P_m=\bra{m}\hat{\rho}_{\rm micro}(E)\ket{m}=\frac{1}{\Omega}\sum_n\delta(E-E_n)\bra{m}\ket{n}\bra{n}\ket{m}=\frac{\delta(E-E_m)}{\Omega}\]
由于$E$只能在$\{E_n\}$上取值，因此对每个态$\ket{m}$，系统处于其上的概率$P_m=1/\Omega$，即能量$E$对应的本征态上每个态的概率权重都相等，满足等概率原理。进而我们可以计算微正则系综的熵：
\[S=-k_{\rm B}\Tr[\hat{\rho}_{\rm micro}\ln\hat{\rho}_{\rm micro}]=-k_{\rm B}\sum_n\frac{1}{\Omega}\ln\frac{1}{\Omega}=-k_{\rm B}\Omega\cdot\frac{1}{\Omega}\ln\frac{1}{\Omega}=k_{\rm B}\ln\Omega\]
\end{proposition}

\begin{proposition}[量子化的正则系综]
类似的对正则系综，自然变量是$(N,V,T)$，将经典正则系综分布函数和配分函数算符化可得：
\[\hat{\rho}_{\rm can}=\frac{e^{-\beta\hat{H}}}{Z}=\frac{1}{Z}\sum_n e^{-\beta E_n}\ket{n}\bra{n}, \quad Z=\Tr\left[e^{-\beta\hat{H}}\right]=\sum_n e^{-\beta E_n}\]
计算正则系综的内能$U$,熵$S$和 Helmholtz 自由能$F$：
\[\begin{aligned}
U&=\Tr(\hat{\rho}_{\rm can}\hat{H})=\frac{1}{Z}\Tr\left[e^{-\beta\hat{H}}\hat{H}\right]=-\pdv{\beta}\ln\left[\Tr e^{-\beta\hat{H}}\right]=-\pdv{\ln Z}{\beta}\\
S&=-k_{\rm B}\Tr[\hat{\rho}_{\rm can}\ln\hat{\rho}_{\rm can}]=-k_{\rm B}\Tr\left[\hat{\rho}_{\rm can}\left(-\beta\hat{H}-\hat{I}\ln Z\right)\right]=\frac{U}{T}+k_{\rm B}\ln Z\\
F&=U-TS=-k_{\rm B}T\ln Z
\end{aligned}\]
\end{proposition}

\begin{proposition}[量子化的巨正则系综]
最后考虑巨正则系综的自然变量是$(T,V,\mu)$，将经典巨正则系综分布函数和配分函数算符化可得：
\[\hat{\rho}_{\rm gc}=\frac{e^{-\beta(\hat{H}-\mu\hat{N})}}{\Xi}, \quad \Xi=\Tr\left[e^{-\beta(\hat{H}-\mu\hat{N})}\right]=\sum_{n=0}^{\infty}Z(N)e^{\beta\mu N}\]
计算巨正则系综的内能$U$，平均粒子数$\expval{N}$，熵$S$和巨势$\Phi_{\rm G}$：
\[\begin{aligned}
U&=\Tr(\hat{\rho}_{\rm gc}\hat{H})=\frac{1}{\Xi}\Tr\left[e^{-\beta(\hat{H}-\mu\hat{N})}\hat{H}\right]=-\pdv{\beta}\ln\left[\Tr e^{-\beta(\hat{H}-\mu\hat{N})}\right]+\mu \expval{N}=-\pdv{\ln\Xi}{\beta}+\mu \expval{N}\\
\expval{N}&=\Tr(\hat{\rho}_{\rm gc}\hat{N})=\frac{1}{\Xi}\Tr\left[e^{-\beta(\hat{H}-\mu\hat{N})}\hat{N}\right]=\frac{1}{\beta}\pdv{\ln\Xi}{\mu}\\
S&=-k_{\rm B}\Tr[\hat{\rho}_{\rm gc}\ln\hat{\rho}_{\rm gc}]=-k_{\rm B}\Tr\left[\hat{\rho}_{\rm gc}\left(-\beta(\hat{H}-\mu\hat{N})-\hat{I}\ln \Xi\right)\right]=\frac{U}{T}+k_{\rm B}\ln \Xi\\
\Phi_{\rm G}&=U-TS-\mu\expval{N}=-k_{\rm B}T\ln\Xi
\end{aligned}\]
\end{proposition}
上面我们只介绍了三个系综中最重要的统计物理量，其他统计物理量也可以从配分函数出发推导获得。下面我们来看一个具体的例子：考虑温度$T$下体积为$V$的三维空间中单个自由粒子，满足薛定谔方程：
\[\hat{H}\ket{\mathbf{k}}=\frac{\hat{\mathbf{p}}^2}{2m}\ket{\mathbf{k}}=\frac{\hbar^2\mathbf{k}^2}{2m}\ket{\mathbf{k}}=\varepsilon_{\mathbf{k}}\ket{\mathbf{k}}\]
对于体积为$V$的三维盒子，在周期性边界条件下每个$\mathbf{k}$态在$k$空间中占据体积$(2\pi)^3/V$，态$\ket{\mathbf{k}}$在位置表象下的表示为平面波$\bra{\mathbf{r}}\ket{\mathbf{k}}=e^{i\mathbf{k}\cdot\mathbf{r}}/\sqrt{V}$。
在正则系综中，单粒子配分函数为：
\[z_1=\Tr\left[e^{-\beta\hat{H}}\right]=\frac{V}{(2\pi)^3}\int\dd[3]{\mathbf{k}}\exp\left(-\frac{\beta\hbar^2 \mathbf{k}^2}{2m}\right)=\frac{V}{(2\pi)^3}\cdot 4\pi\int_0^{\infty}k^2\dd{k}\exp\left(-\frac{\beta\hbar^2k^2}{2m}\right)=\frac{V}{\Lambda^3}\]
其中$\Lambda=h/\sqrt{2\pi m k_{\rm B} T}$是热波长。容易获得正则分布函数和其在动量表象下的矩阵元：
\[\hat{\rho}_{\rm can}=\frac{e^{-\beta\hat{H}}}{z_1}=\frac{1}{z_1}\sum_{\mathbf{k}}e^{-\beta\varepsilon_{\mathbf{k}}}\ket{\mathbf{k}}\bra{\mathbf{k}} \quad \Rightarrow \quad \bra{\mathbf{k}}\hat{\rho}_{\rm can}\ket{\mathbf{k}'}=\frac{1}{z_1}\exp\left(-\frac{\beta\hbar^2\mathbf{k}^2}{2m}\right)\delta_{\mathbf{k}\mathbf{k}'}\]
更进一步也可以获得正则分布函数在位置表象下的矩阵元，具体的推导见 [章节\ref{sec:gaussian}]：
\[\begin{aligned}
\bra{\mathbf{r}}\hat{\rho}_{\rm can}\ket{\mathbf{r}'}&=\frac{1}{z_1}\sum_{\mathbf{k}}e^{-\beta\varepsilon_{\mathbf{k}}}\bra{\mathbf{r}}\ket{\mathbf{k}}\bra{\mathbf{k}}\ket{\mathbf{r}'}=\frac{1}{z_1V}\sum_{\mathbf{k}}e^{-\beta\varepsilon_{\mathbf{k}}}e^{i\mathbf{k}\cdot(\mathbf{r}-\mathbf{r}')}=\frac{1}{z_1V}\int\frac{V}{(2\pi)^3}\dd[3]{\mathbf{k}} e^{-\beta\varepsilon_{\mathbf{k}}}e^{i\mathbf{k}\cdot(\mathbf{r}-\mathbf{r}')}\\
&=\frac{1}{z_1}\left(\frac{m}{2\pi\beta\hbar^2}\right)^{3/2}\exp\left[-\frac{m|\mathbf{r}-\mathbf{r}'|^2}{2\beta\hbar^2}\right]=\frac{1}{V}\exp\left[-\frac{m|\mathbf{r}-\mathbf{r}'|^2}{2\beta\hbar^2}\right]
\end{aligned}\]

在此基础上我们可以计算内能$U$：
\[\begin{aligned}
U&=\Tr\left[\hat{\rho}_{\rm can}\hat{H}\right]=\sum_{\mathbf{k}}\bra{\mathbf{k}}\hat{\rho}_{\rm can}\hat{H}\ket{\mathbf{k}}=\sum_{\mathbf{k}}\sum_{\mathbf{k}'}\bra{\mathbf{k}}\hat{\rho}_{\rm can}\ket{\mathbf{k}'}\bra{\mathbf{k}'}\hat{H}\ket{\mathbf{k}}\\
&=\sum_{\mathbf{k}}\sum_{\mathbf{k}'}\frac{1}{z_1}\exp\left(-\frac{\beta\hbar^2\mathbf{k}^2}{2m}\right)\delta_{\mathbf{k}\mathbf{k}'}\frac{\hbar^2\mathbf{k}^2}{2m}\delta_{\mathbf{k}\mathbf{k}'}=\sum_{\mathbf{k}}\frac{1}{z_1}\exp\left(-\frac{\beta\hbar^2\mathbf{k}^2}{2m}\right)\frac{\hbar^2\mathbf{k}^2}{2m}\\
&=\frac{1}{z_1}\frac{V}{(2\pi)^3}\int\dd[3]{\mathbf{k}}\exp\left(-\frac{\beta\hbar^2\mathbf{k}^2}{2m}\right)\frac{\hbar^2\mathbf{k}^2}{2m}=\frac{\hbar^2\Lambda^3}{2m(2\pi)^3}\int\mathbf{k}^2\exp\left(-\frac{\beta\hbar^2\mathbf{k}^2}{2m}\right)\dd[3]{\mathbf{k}}\\
&=\frac{\hbar^2\Lambda^3}{2m(2\pi)^3}\cdot 4\pi\int_0^{\infty} k^4\exp\left(-\frac{\beta\hbar^2k^2}{2m}\right)\dd{k}=\frac{\hbar^2\Lambda^3}{m(2\pi)^2}\cdot \left(\frac{2m}{\beta\hbar^2}\right)^{\frac{5}{2}}\int_0^{\infty} x^4e^{-x^2}\dd{x}\\
&=\frac{\hbar^2\Lambda^3}{m(2\pi)^2}\cdot \left(\frac{2m}{\beta\hbar^2}\right)^{\frac{5}{2}}\cdot\frac{1}{2}\int_0^{\infty} t^{\frac{3}{2}}e^{-t}\dd{t}=\frac{\hbar^2\Lambda^3}{m(2\pi)^2}\cdot\frac{1}{2}\left(\frac{2m}{\beta\hbar^2}\right)^{\frac{5}{2}}\Gamma(\frac{5}{2})=\frac{3}{2}\cdot\frac{\Lambda^3}{\beta}\left(\frac{m}{2\pi\beta\hbar^2}\right)^{\frac{3}{2}}=\frac{3}{2}k_{\rm B}T
\end{aligned}\]
或者我们简单一点，从配分函数出发：
\[U=-\pdv{\ln z_1}{\beta}=\frac{3}{2}k_{\rm B}T\]
最后我们来看看热波长，当温度$T\to\infty$时，$\Lambda\to 0$，此时单粒子配分函数$z_1=V/\Lambda^3\to\infty$，系统的状态数变得非常大，量子效应变得微弱；当温度$T\to 0$时，$\Lambda\to\infty$，单粒子配分函数趋于零，系统的状态数变得非常有限，量子效应变得显著。

\newpage
\subsection{理想量子气体与 Bose-Einstein / Fermi-Dirac 分布}
对无相互作用体系，哈密顿量与粒子数算符都可写成占据数之和
\[
\hat{H}=\sum_l \varepsilon_l \hat{n}_l,\qquad \hat{N}=\sum_l \hat{n}_l.
\]
这时巨正则配分函数可以在每个单粒子能级上完全因子化，因此量子理想气体最自然的处理方式就是巨正则系综。

对玻色子，$n_l=0,1,2,\cdots$，于是
\[
\Xi_{\rm BE}
=\prod_l\sum_{n_l=0}^{\infty}e^{-\beta n_l(\varepsilon_l-\mu)}
=\prod_l\frac{1}{1-e^{-\beta(\varepsilon_l-\mu)}}.
\]
对应的巨势为
\[
\Omega_{\rm BE}=k_B T\sum_l\ln\left(1-e^{-\beta(\varepsilon_l-\mu)}\right),
\]
平均占据数则是
\[
\boxed{\bar n_l^{\rm (BE)}=\frac{1}{e^{\beta(\varepsilon_l-\mu)}-1}}.
\]

对费米子，$n_l=0,1$，因此
\[
\Xi_{\rm FD}
=\prod_l\sum_{n_l=0}^{1}e^{-\beta n_l(\varepsilon_l-\mu)}
=\prod_l\left(1+e^{-\beta(\varepsilon_l-\mu)}\right),
\]
对应巨势
\[
\Omega_{\rm FD}=-k_B T\sum_l\ln\left(1+e^{-\beta(\varepsilon_l-\mu)}\right),
\]
以及平均占据数
\[
\boxed{\bar n_l^{\rm (FD)}=\frac{1}{e^{\beta(\varepsilon_l-\mu)}+1}}.
\]
这两个分布在形式上极其相似，但分母中的符号差异反映了截然不同的量子统计：玻色子倾向于“扎堆”占据同一个态，而费米子则受到 Pauli 不相容原理约束，每个单粒子态最多只能放一个粒子。

从占据数分布出发，粒子总数和内能都可以写成单粒子态求和：
\[
\boxed{\bar N=\sum_l \bar n_l},\qquad 
\boxed{U=\sum_l \varepsilon_l \bar n_l}.
\]
单个能级上的涨落也可直接算出：
\[
\expval{(\Delta n_l)^2}_{\rm BE}=\bar n_l(1+\bar n_l),\qquad
\expval{(\Delta n_l)^2}_{\rm FD}=\bar n_l(1-\bar n_l).
\]
玻色子的正号表明占据存在 bunching 倾向，而费米子的负号则对应反聚束与占据抑制。

在高温低密度极限下，若所有能级都满足$\bar n_l\ll 1$，则两种量子分布统一退化为
\[
\bar n_l \approx e^{-\beta(\varepsilon_l-\mu)},
\]
这就是 Maxwell-Boltzmann 分布。因此经典统计力学并不是量子统计力学之外的另一套理论，而是量子统计在稀薄极限下的近似。

对自由粒子，若单粒子能量为$\varepsilon_{\mathbf{p}}=p^2/2m$，并把求和改为动量积分，则可以进一步得到量子理想气体的状态方程。对玻色子必须满足$\mu<\varepsilon_0$，在低温高密度下会出现基态宏观占据，这就是 Bose-Einstein 凝聚；而对费米子，在$T\to 0$时所有$\varepsilon_l<\mu$的能级被逐个填满，形成费米面。这两种现象都是纯粹由量子统计导致的，在经典极限中都不会出现。

\subsection{两个典型例子：自由自旋与量子谐振子}
作为例子，先看处在外磁场$B$中的$N$个彼此独立的自旋$\frac{1}{2}$粒子。设单个自旋的磁矩大小为$\mu$，则其两能级为
\[
\varepsilon_{\pm}=\mp \mu B.
\]
单个粒子的配分函数为
\[
z_1=e^{\beta\mu B}+e^{-\beta\mu B}=2\cosh(\beta\mu B),
\]
于是$N$个独立自旋的正则配分函数就是
\[
Z_N=z_1^N=\left[2\cosh(\beta\mu B)\right]^N.
\]
由此立刻得到自由能和内能
\[
F=-Nk_B T\ln\left[2\cosh(\beta\mu B)\right],
\]
\[
U=-\pdv{\ln Z_N}{\beta}=-N\mu B\tanh(\beta\mu B).
\]
磁化强度可由自由能对磁场求导获得：
\[
M=-\left(\pdv{F}{B}\right)_{T,N}=N\mu\tanh(\beta\mu B).
\]
相应的热容为
\[
C_V=Nk_B(\beta\mu B)^2\,\text{sech}^2(\beta\mu B).
\]
这个模型虽然简单，却已经展示了量子能级离散化如何直接决定热容、磁化率和涨落行为。

再看频率为$\omega$的一维量子谐振子，其能级为
\[
\varepsilon_n=\hbar\omega\left(n+\frac{1}{2}\right),\qquad n=0,1,2,\ldots.
\]
单个谐振子的配分函数可以直接求和得到
\[
z_1=\sum_{n=0}^{\infty}e^{-\beta\hbar\omega(n+1/2)}
=\frac{e^{-\beta\hbar\omega/2}}{1-e^{-\beta\hbar\omega}}.
\]
于是
\[
F=-k_B T\ln z_1
=\frac{1}{2}\hbar\omega+k_B T\ln\left(1-e^{-\beta\hbar\omega}\right),
\]
\[
U=-\pdv{\ln z_1}{\beta}
=\hbar\omega\left(\frac{1}{2}+\frac{1}{e^{\beta\hbar\omega}-1}\right).
\]
进一步得到热容
\[
C_V=k_B\frac{(\beta\hbar\omega)^2e^{\beta\hbar\omega}}{(e^{\beta\hbar\omega}-1)^2}.
\]
在高温极限$\beta\hbar\omega\ll 1$时，上式趋于$k_B$，回到经典等分配定理；而在低温极限下，热容指数衰减到$0$，这正是量子化能级结构的直接后果。

综上所述，量子统计力学只是把“经典相空间中的概率密度”替换成了“希尔伯特空间中的密度算符”，但这一替换会带来两个深刻的新现象：一是熵和热力学量必须通过迹运算来定义；二是全同粒子的交换对称性会把系统分成玻色子与费米子两大类，并最终导致 Bose-Einstein 与 Fermi-Dirac 统计。理解了这两点，后续无论是处理电子气、光子气、声子、磁性模型还是凝聚态多体系统，都会有统一的出发点。
